\section{完整案例：前后对比与目标评价}
下面用同一尺度比较方法改进前后的得分，并将目标偏差单独表达。
数据为演示值；哑铃图例写在表注，避免读者混淆两个端点。
\begin{example}[text and listing,sidebyside=false]{示例 71：实验前后对比表}
\gradbarsstyle{study}{preset=paper,width=48mm,
  min=0,max=100,precision=0,palette=categorical}
\gradbarsstyle{study-change}{style=study,
  track color=black!40,stem width=.9pt}
\begin{center}
\begin{tabular}{@{}ccc@{}}
\toprule
\multicolumn{1}{c}{Method} & \multicolumn{1}{c}{After / Before} & \multicolumn{1}{c}{Target 80}\\\midrule
A & \graddumbbell[style=study-change]{58}{76}
  & \gradbar[style=study,width=23mm,better=target,
    quality target=80,thresholds={5,15}]{76}\\
B & \graddumbbell[style=study-change]{63}{84}
  & \gradbar[style=study,width=23mm,better=target,
    quality target=80,thresholds={5,15}]{84}\\
C & \graddumbbell[style=study-change]{65}{65}
  & \gradbar[style=study,width=23mm,better=target,
    quality target=80,thresholds={5,15}]{65}\\
\bottomrule
\end{tabular}
\end{center}
\end{example}
空心圆为 Before，实心圆为 After；目标评价条越长仍表示数值越大。
颜色表示到 80 的偏差，而非“越长越好”。

背景轨道采用 \key{track color=black!40} 加深，并以 \key{stem width=.9pt} 增强连线；
灰线表示完整的 0--100 量程，蓝线仅连接前后两个观测值。
